\documentclass[a4paper,12pt,twoside]{memoir}

% Castellano
\usepackage[spanish,es-tabla]{babel}
\selectlanguage{spanish}
\usepackage[utf8]{inputenc}
\usepackage[T1]{fontenc}
\usepackage{lmodern} % Scalable font
\usepackage{microtype}
\usepackage{placeins}

\RequirePackage{booktabs}
\RequirePackage[table]{xcolor}
\RequirePackage{xtab}
\RequirePackage{multirow}

% Links
\PassOptionsToPackage{hyphens}{url}\usepackage[colorlinks]{hyperref}
\hypersetup{
	allcolors = {red}
}

% Ecuaciones
\usepackage{amsmath}

% Rutas de fichero / paquete
\newcommand{\ruta}[1]{{\sffamily #1}}

% Párrafos
\nonzeroparskip

% Huérfanas y viudas
\widowpenalty100000
\clubpenalty100000

% Imagenes
\usepackage{graphicx}
\newcommand{\imagen}[2]{
	\begin{figure}[!h]
		\centering
		\includegraphics[width=0.9\textwidth]{#1}
		\caption{#2}\label{fig:#1}
	\end{figure}
	\FloatBarrier
}

\newcommand{\imagenflotante}[2]{
	\begin{figure}%[!h]
		\centering
		\includegraphics[width=0.9\textwidth]{#1}
		\caption{#2}\label{fig:#1}
	\end{figure}
}



% El comando \figura nos permite insertar figuras comodamente, y utilizando
% siempre el mismo formato. Los parametros son:
% 1 -> Porcentaje del ancho de página que ocupará la figura (de 0 a 1)
% 2 --> Fichero de la imagen
% 3 --> Texto a pie de imagen
% 4 --> Etiqueta (label) para referencias
% 5 --> Opciones que queramos pasarle al \includegraphics
% 6 --> Opciones de posicionamiento a pasarle a \begin{figure}
\newcommand{\figuraConPosicion}[6]{%
  \setlength{\anchoFloat}{#1\textwidth}%
  \addtolength{\anchoFloat}{-4\fboxsep}%
  \setlength{\anchoFigura}{\anchoFloat}%
  \begin{figure}[#6]
    \begin{center}%
      \Ovalbox{%
        \begin{minipage}{\anchoFloat}%
          \begin{center}%
            \includegraphics[width=\anchoFigura,#5]{#2}%
            \caption{#3}%
            \label{#4}%
          \end{center}%
        \end{minipage}
      }%
    \end{center}%
  \end{figure}%
}

%
% Comando para incluir imágenes en formato apaisado (sin marco).
\newcommand{\figuraApaisadaSinMarco}[5]{%
  \begin{figure}%
    \begin{center}%
    \includegraphics[angle=90,height=#1\textheight,#5]{#2}%
    \caption{#3}%
    \label{#4}%
    \end{center}%
  \end{figure}%
}
% Para las tablas
\newcommand{\otoprule}{\midrule [\heavyrulewidth]}
%
% Nuevo comando para tablas pequeñas (menos de una página).
\newcommand{\tablaSmall}[5]{%
 \begin{table}
  \begin{center}
   \rowcolors {2}{gray!35}{}
   \begin{tabular}{#2}
    \toprule
    #4
    \otoprule
    #5
    \bottomrule
   \end{tabular}
   \caption{#1}
   \label{tabla:#3}
  \end{center}
 \end{table}
}

%
% Nuevo comando para tablas pequeñas (menos de una página).
\newcommand{\tablaSmallSinColores}[5]{%
 \begin{table}[H]
  \begin{center}
   \begin{tabular}{#2}
    \toprule
    #4
    \otoprule
    #5
    \bottomrule
   \end{tabular}
   \caption{#1}
   \label{tabla:#3}
  \end{center}
 \end{table}
}

\newcommand{\tablaApaisadaSmall}[5]{%
\begin{landscape}
  \begin{table}
   \begin{center}
    \rowcolors {2}{gray!35}{}
    \begin{tabular}{#2}
     \toprule
     #4
     \otoprule
     #5
     \bottomrule
    \end{tabular}
    \caption{#1}
    \label{tabla:#3}
   \end{center}
  \end{table}
\end{landscape}
}

%
% Nuevo comando para tablas grandes con cabecera y filas alternas coloreadas en gris.
\newcommand{\tabla}[6]{%
  \begin{center}
    \tablefirsthead{
      \toprule
      #5
      \otoprule
    }
    \tablehead{
      \multicolumn{#3}{l}{\small\sl continúa desde la página anterior}\\
      \toprule
      #5
      \otoprule
    }
    \tabletail{
      \hline
      \multicolumn{#3}{r}{\small\sl continúa en la página siguiente}\\
    }
    \tablelasttail{
      \hline
    }
    \bottomcaption{#1}
    \rowcolors {2}{gray!35}{}
    \begin{xtabular}{#2}
      #6
      \bottomrule
    \end{xtabular}
    \label{tabla:#4}
  \end{center}
}

%
% Nuevo comando para tablas grandes con cabecera.
\newcommand{\tablaSinColores}[6]{%
  \begin{center}
    \tablefirsthead{
      \toprule
      #5
      \otoprule
    }
    \tablehead{
      \multicolumn{#3}{l}{\small\sl continúa desde la página anterior}\\
      \toprule
      #5
      \otoprule
    }
    \tabletail{
      \hline
      \multicolumn{#3}{r}{\small\sl continúa en la página siguiente}\\
    }
    \tablelasttail{
      \hline
    }
    \bottomcaption{#1}
    \begin{xtabular}{#2}
      #6
      \bottomrule
    \end{xtabular}
    \label{tabla:#4}
  \end{center}
}

%
% Nuevo comando para tablas grandes sin cabecera.
\newcommand{\tablaSinCabecera}[5]{%
  \begin{center}
    \tablefirsthead{
      \toprule
    }
    \tablehead{
      \multicolumn{#3}{l}{\small\sl continúa desde la página anterior}\\
      \hline
    }
    \tabletail{
      \hline
      \multicolumn{#3}{r}{\small\sl continúa en la página siguiente}\\
    }
    \tablelasttail{
      \hline
    }
    \bottomcaption{#1}
  \begin{xtabular}{#2}
    #5
   \bottomrule
  \end{xtabular}
  \label{tabla:#4}
  \end{center}
}



\definecolor{cgoLight}{HTML}{EEEEEE}
\definecolor{cgoExtralight}{HTML}{FFFFFF}

%
% Nuevo comando para tablas grandes sin cabecera.
\newcommand{\tablaSinCabeceraConBandas}[5]{%
  \begin{center}
    \tablefirsthead{
      \toprule
    }
    \tablehead{
      \multicolumn{#3}{l}{\small\sl continúa desde la página anterior}\\
      \hline
    }
    \tabletail{
      \hline
      \multicolumn{#3}{r}{\small\sl continúa en la página siguiente}\\
    }
    \tablelasttail{
      \hline
    }
    \bottomcaption{#1}
    \rowcolors[]{1}{cgoExtralight}{cgoLight}

  \begin{xtabular}{#2}
    #5
   \bottomrule
  \end{xtabular}
  \label{tabla:#4}
  \end{center}
}



\graphicspath{ {./img/} }

% Capítulos
\chapterstyle{bianchi}
\newcommand{\capitulo}[2]{
	\setcounter{chapter}{#1}
	\setcounter{section}{0}
	\setcounter{figure}{0}
	\setcounter{table}{0}
	\chapter*{#2}
	\addcontentsline{toc}{chapter}{#2}
	\markboth{#2}{#2}
}

% Apéndices
\renewcommand{\appendixname}{Apéndice}
\renewcommand*\cftappendixname{\appendixname}

\newcommand{\apendice}[1]{
	%\renewcommand{\thechapter}{A}
	\chapter{#1}
}

\renewcommand*\cftappendixname{\appendixname\ }

% Formato de portada
\makeatletter
\usepackage{xcolor}
\newcommand{\tutor}[1]{\def\@tutor{#1}}
\newcommand{\course}[1]{\def\@course{#1}}
\definecolor{cpardoBox}{HTML}{E6E6FF}
\def\maketitle{
  \null
  \thispagestyle{empty}
  % Cabecera ----------------
\noindent\includegraphics[width=\textwidth]{cabecera}\vspace{1cm}%
  \vfill
  % Título proyecto y escudo informática ----------------
  \colorbox{cpardoBox}{%
    \begin{minipage}{.8\textwidth}
      \vspace{.5cm}\Large
      \begin{center}
      \textbf{TFG del Grado en Ingeniería Informática}\vspace{.6cm}\\
      \textbf{\LARGE\@title{}}
      \end{center}
      \vspace{.2cm}
    \end{minipage}

  }%
  \hfill\begin{minipage}{.20\textwidth}
    \includegraphics[width=\textwidth]{escudoInfor}
  \end{minipage}
  \vfill
  % Datos de alumno, curso y tutores ------------------
  \begin{center}%
  {%
    \noindent\LARGE
    Presentado por \@author{}\\ 
    en Universidad de Burgos --- \@date{}\\
    Tutores: \@tutor{}\\
  }%
  \end{center}%
  \null
  \cleardoublepage
  }
\makeatother

\newcommand{\nombre}{Nombre del alumno} %%% cambio de comando

% Datos de portada
\title{Aplicación Web para Gestión de Datos Genómicos de Cánceres de Pulmón}
\author{Carlos Ortúñez Rojo}
\tutor{Dr. Jesús Manuel Maudes Raedo y Dr. José Francisco Díez Pastor}
\date{\today}

\begin{document}

\maketitle


\newpage\null\thispagestyle{empty}\newpage


%%%%%%%%%%%%%%%%%%%%%%%%%%%%%%%%%%%%%%%%%%%%%%%%%%%%%%%%%%%%%%%%%%%%%%%%%%%%%%%%%%%%%%%%
\thispagestyle{empty}


\noindent\includegraphics[width=\textwidth]{cabecera}\vspace{1cm}

\noindent D. nombre tutor, profesor del departamento de nombre departamento, área de nombre área.

\noindent Expone:

\noindent Que el alumno D. \nombre, con DNI dni, ha realizado el Trabajo final de Grado en Ingeniería Informática titulado título de TFG. 

\noindent Y que dicho trabajo ha sido realizado por el alumno bajo la dirección del que suscribe, en virtud de lo cual se autoriza su presentación y defensa.

\begin{center} %\large
En Burgos, {\large \today}
\end{center}

\vfill\vfill\vfill

% Author and supervisor
\begin{minipage}{0.45\textwidth}
\begin{flushleft} %\large
Vº. Bº. del Tutor:\\[2cm]
D. nombre tutor
\end{flushleft}
\end{minipage}
\hfill
\begin{minipage}{0.45\textwidth}
\begin{flushleft} %\large
Vº. Bº. del co-tutor:\\[2cm]
D. nombre co-tutor
\end{flushleft}
\end{minipage}
\hfill

\vfill

% para casos con solo un tutor comentar lo anterior
% y descomentar lo siguiente
%Vº. Bº. del Tutor:\\[2cm]
%D. nombre tutor


\newpage\null\thispagestyle{empty}\newpage




\frontmatter

% Abstract en castellano
\renewcommand*\abstractname{Resumen}
\begin{abstract}
En este proyecto se ha a desarrollado una aplicación web que permite la introducción, modificación, eliminación y visualización de los datos genómicos de canceres de pulmón de los pacientes del Hospital Universitario de Burgos. Esta aplicación web también cuenta con un sistema de autenticación y un sistema de gestión de usuarios.

Todos los datos están almacenados en una base de datos relacional para permitir un rápido acceso a estos y ahorrar espacio a la hora de almacenarlos.

Dentro de la visualización de los datos se han incluido varias opciones para permitir a los oncólogos una mejor explotación a la hora de realizar sus estudios. Estas opciones de visualización incluyen la exportación a un fichero \textit{.xlsx}, la representación en diagramas tanto de sectores como de columnas, la obtención de tablas de frecuencias y de percentiles y la visualización individual de los datos crudos de cada paciente.

Para poder comprobar el correcto funcionamiento de toda la parte de visualización nombrada anteriormente, se ha realizado en \textit{Python} un script para poder generar una base de datos sintética permitiendo al usuario elegir diferentes parámetros de las distribuciones que se han considerado más óptimas para el tipo de datos con los que se está trabajando.
\end{abstract}

\renewcommand*\abstractname{Descriptores}
\begin{abstract}
Aplicación web, base de datos relacional, Laravel, estudio estadístico, visualización, datos médicos.
\end{abstract}

\clearpage

% Abstract en inglés
\renewcommand*\abstractname{Abstract}
\begin{abstract}
This project consists in the development of a web application that allows the input, modification, elimination and visualization of the genomic data regarding lung cancers cancer patients from the Hospital Universitario de Burgos. In addition, this web application will have an authentication system and a user management system.

All the data is stored in a relational database to allow quick access and to save space when storing the information.

Within the data visualization, several options have been included to offer oncologists a better exploitation when conducting their studies. These display options include exporting to a file \textit{.xlsx}, the representation in diagrams of both, sectors and columns, the obtention of tables than show frequencies and percentiles and the individual visualization of each patient’s raw data.

In order to check the correct performance of the entire display part named above, a \textit{Python} script has been created to be able to generate a synthetic database allowing the user to choose different parameters of the distributions that have been considered most optimal for the type of data we are working with.
\end{abstract}

\renewcommand*\abstractname{Keywords}
\begin{abstract}
Web application, relational database, Laravel, statistical study, visualization, medical data.
\end{abstract}

\clearpage

% Indices
\tableofcontents

\clearpage

\listoffigures

\clearpage

\listoftables
\clearpage

\mainmatter
\capitulo{1}{Introducción}
El \textbf{cáncer} es un conjunto de enfermedades relacionadas las cuales son provocadas por una anomalía en las células que ocasiona que estas empiecen a dividirse continuamente y provoquen un mal funcionamiento de la zona incluso llegándose a expandir a otras partes del cuerpo, a esto se le llama \textbf{metástasis} \cite{cancerweb:online}. El \textbf{cáncer de pulmón} ocurre cuando este problema surge en el tejido pulmonar, este cáncer es el que más muertes provoca actualmente en el mundo, tanto en hombres como en mujeres, siendo el tabaquismo la principal causa de este problema llegando a suponer hasta el 80\% de los casos \cite{cancerpulmon:online}.

Existe gran variedad en cuanto a los diferentes tipos de \textbf{cáncer de pulmón}, los \textbf{tratamientos} que se llevan a cabo y las \textbf{herramientas de detección} de las que se hacen uso. Dicha pluralidad hace que pueda resultar complicada la toma de decisiones durante el proceso de tratamiento del paciente. Además, tratar eficazmente el problema y adaptarse a las circunstancias personales de cada uno de ellos, puede suponer alargar la vida o incluso salvar a la persona. Por ello, realizar una gestión y organización de los datos eficaz, y contar con herramientas que permitan un fácil acceso y visualización, pueden resultar fundamentales en la toma de decisiones de los tratamientos.


En la actualidad el \textbf{HUBU} almacena los datos de los pacientes que han sufrido o sufren cáncer de pulmón en formato texto. Este sistema de visualización conlleva a un análisis poco eficaz de los datos, ya que supone un proceso lento y costoso.

Para solventar este problema, el objetivo de este proyecto ha sido crear una \textbf{aplicación web} que sirva como \textbf{interfaz} para gestionar todos los datos de los pacientes. Estos estarán almacenados dentro de una \textbf{base de datos relacional} y la web tendrá herramientas de \textbf{visualización} de los datos como pueden ser gráficas o tablas permitiendo así, un óptimo análisis de dichos datos y facilitando la toma de decisiones durante los procesos de tratamiento.

También debido a que es posible que los médicos que trabajen con esta  \textbf{aplicación web} no tengan un alto nivel de informática se ha intentado realizar de la manera más \textbf{intuitiva} posible, con muchos  \textbf{controles de errores} y con una comunicación constante con el equipo de oncólogos del Hospital Universitario de Burgos para lograr la mejor versión posible.

\section{Material ajunto}

En el proyecto como material adjunto se incluye: 
\begin{itemize}
	\item Memoria
	\item Anexos 
	\item Aplicación web en \textit{Laravel}.
	\item Scripts \textit{SQL} para generar la base de datos e insertar datos necesarios para la aplicación.
	\item Script en \textit{Python} para generar la base de datos sintética.
\end{itemize}

Todos los elementos desarrollados durante el proyecto se pueden ver en el repositorio de \textit{GitHub}: \url{https://github.com/CarlosOrtu/TFG_HUBU}


\capitulo{2}{Objetivos del proyecto}

En este apartado se van a especificar los diferentes objetivos, tanto funcionales como técnicos.

\section{Objetivos funcionales}

Los objetivos funcionales que se han tratado de conseguir con este proyecto son:

\begin{itemize}
	\item Memoria
	\item Anexos 
	\item Aplicación web en \textit{Laravel}.
	\item Scripts \textit{SQL} para generar la base de datos e insertar datos necesarios para la aplicación.
	\item Script en \textit{Python} para generar la base de datos sintética.
\end{itemize}
\section{Objetivos técnicos}
\capitulo{3}{Conceptos teóricos}

En aquellos proyectos que necesiten para su comprensión y desarrollo de unos conceptos teóricos de una determinada materia o de un determinado dominio de conocimiento, debe existir un apartado que sintetice dichos conceptos.

Algunos conceptos teóricos de \LaTeX \footnote{Créditos a los proyectos de Álvaro López Cantero: Configurador de Presupuestos y Roberto Izquierdo Amo: PLQuiz}.

\section{Secciones}

Las secciones se incluyen con el comando section.

\subsection{Subsecciones}

Además de secciones tenemos subsecciones.

\subsubsection{Subsubsecciones}

Y subsecciones. 


\section{Referencias}

Las referencias se incluyen en el texto usando cite \cite{wiki:latex}. Para citar webs, artículos o libros \cite{koza92}.


\section{Imágenes}

Se pueden incluir imágenes con los comandos standard de \LaTeX, pero esta plantilla dispone de comandos propios como por ejemplo el siguiente:

\imagen{escudoInfor}{Autómata para una expresión vacía}



\section{Listas de items}

Existen tres posibilidades:

\begin{itemize}
	\item primer item.
	\item segundo item.
\end{itemize}

\begin{enumerate}
	\item primer item.
	\item segundo item.
\end{enumerate}

\begin{description}
	\item[Primer item] más información sobre el primer item.
	\item[Segundo item] más información sobre el segundo item.
\end{description}
	
\begin{itemize}
\item 
\end{itemize}

\section{Tablas}

Igualmente se pueden usar los comandos específicos de \LaTeX o bien usar alguno de los comandos de la plantilla.

\tablaSmall{Herramientas y tecnologías utilizadas en cada parte del proyecto}{l c c c c}{herramientasportipodeuso}
{ \multicolumn{1}{l}{Herramientas} & App AngularJS & API REST & BD & Memoria \\}{ 
HTML5 & X & & &\\
CSS3 & X & & &\\
BOOTSTRAP & X & & &\\
JavaScript & X & & &\\
AngularJS & X & & &\\
Bower & X & & &\\
PHP & & X & &\\
Karma + Jasmine & X & & &\\
Slim framework & & X & &\\
Idiorm & & X & &\\
Composer & & X & &\\
JSON & X & X & &\\
PhpStorm & X & X & &\\
MySQL & & & X &\\
PhpMyAdmin & & & X &\\
Git + BitBucket & X & X & X & X\\
Mik\TeX{} & & & & X\\
\TeX{}Maker & & & & X\\
Astah & & & & X\\
Balsamiq Mockups & X & & &\\
VersionOne & X & X & X & X\\
} 

\capitulo{4}{Técnicas y herramientas}

\section{Técnicas}

En este apartado se va a explicar la técnica de desarrollo de software que ha seguido y una explicación de porque se ha optado por esta opción y como se ha adaptado al proyecto.

\subsection{Scrum}

Las metodología ágiles son aquellas que se basan en dar una mayor libertad a la colaboración con el cliente y a la realización del proyecto en iteraciones. Según el manifiesto ágil esta nueva metodología se basa en cuatro principios: Priorizar al individuo y a las iteraciones sobre el proceso y las herramientas, el desarrollo del software sobre la documentación, colaboración con el cliente antes que un contrato fijo y la capacidad de responder a los cambios antes que un plan fijo\cite{metodologia_agil}.

Debido a las características del proyecto se decidió que una metodología ágil era lo más adecuado para lleva a cabo su desarrollo ya que se comenzó con una idea general, por lo tanto a la hora de ir realizando el proyecto iba a ser necesario realizar cambios sobre las partes ya finalizadas y en vista de que a medida que avanzará el proyecto se iban a ir añadiendo funcionalidades, por lo tanto no se puede establecer un plan fijo desde el comienzo\cite{metodologia_agil}.

De todas las metodologías ágiles para desarrollar el proyecto se eligió Scrum. Este método se basa en el desarrollo de sprints, en los cuales ser realizan partes del proyecto final ya funcionales y reuniones periódicas. Estos sprints suelen durar entre 1 ó 2 semanas\cite{scrum}. Esta metodología está pensada para trabajar con equipos enteros de desarrollo dentro de una empresa, pero para este proyecto se ha adaptado. Los roles quedan divididos de la manera que el \textit{product owner} son los oncólogos, el \textit{scrumMaster} son ambos tutores y el equipo de desarrollo soy yo, las reuniones y las duraciones de los sprint también han sido levemente modificadas haciéndolas menos estrictas debido a que no siempre era posible para los oncólogos realizar una reunión a la finalización de los sprint.

\section{Herramientas} 

En este apartado se van a explicar todas las herramientas usadas para llevar a cabo cada parte del proyecto y se van a nombrar las diferentes opciones que se tuvieron en cuenta .

\subsection{Infraestructura de desarrollo local}

\subsubsection{WAMP} \label{WAMP}
\begin{itemize}
    \item Herramientas consideradas: \href{https://www.mamp.info/en/windows/}{MAMP}\footnote{\href{https://www.mamp.info/en/windows/}{https://www.mamp.info/en/windows/}}, \href{https://www.apachefriends.org/es/index.html}{XAMPP}\footnote{\href{https://www.apachefriends.org/es/index.html}{https://www.apachefriends.org/es/index.html}}, \href{https://www.wampserver.com/en/}{WAMP}\footnote{\label{wampfoot}\href{https://www.wampserver.com/en/}{https://www.wampserver.com/en/}}.
	\item Herramienta elegida: \href{https://www.wampserver.com/en/}{WAMP}\footref{wampfoot}.
\end{itemize}

\href{http://www.wampserver.es/}{WAMP}\footref{wampfoot} es un entorno de desarrollo web local pensando para usarse en Windows el cual contiene Apache, MySQL, PHP y el administrador de bases de datos PhpMyAdmin. Por lo tanto este entorno cumple con los objetivos técnicos que se tenían para el desarrollo de la aplicación web.

\subsection{Host web}

\subsubsection{Heroku} \label{Heroku}
\begin{itemize}
    \item Herramientas consideradas: \href{https://www.ionos.es/alojamiento/alojamiento-web}{Ionos}\footnote{\href{https://www.ionos.es/alojamiento/alojamiento-web}{https://www.ionos.es/alojamiento/alojamiento-web}}, \href{https://aws.amazon.com/es/ec2/dedicated-hosts/}{AWS amazon}\footnote{\href{https://aws.amazon.com/es/ec2/dedicated-hosts/}{https://aws.amazon.com/es/ec2/dedicated-hosts/}}, \href{https://www.heroku.com/}{Heroku}\footnote{\label{herokufoot}\href{https://www.heroku.com/}{https://www.heroku.com/}}.
    \item Herramienta elegida: \href{https://www.heroku.com/}{Heroku}\footref{herokufoot}.
\end{itemize}

\href{https://www.heroku.com/}{Heroku}\footref{herokufoot} es una plataforma de computación en la nube que da soporte a varios lenguajes como Python, Java y PHP\cite{heroku}. También cuenta con una serie de extensiones que resultaban interesantes para nuestro proyecto, estas son, una base de datos MySQL que permite la gestión desde MySQL Workbench y un protocolo HTTPS para implementar a nuestra web. También para aplicaciones pequeñas como la que se esta desarrollando en este trabajo el host resulta gratuito y la política de privacidad de datos concordaba con los requisitos necesarios para poder alojar los datos de los pacientes.

\subsection{Back-end}

\subsubsection{Laravel}\label{Laravel}
\begin{itemize}
    \item Herramientas consideradas: \href{https://www.djangoproject.com/}{django}\footnote{\href{https://www.djangoproject.com/}{https://www.djangoproject.com/}}, \href{https://docs.microsoft.com/es-es/dotnet/}{.NET}\footnote{\href{https://docs.microsoft.com/es-es/dotnet/}{https://docs.microsoft.com/es-es/dotnet/}}, \href{https://laravel.com/}{Laravel}\footnote{\label{laravelfoot}\href{https://laravel.com/}{https://laravel.com/}}.
	\item Herramienta elegida: \href{https://laravel.com/}{Laravel}\footref{laravelfoot}.
\end{itemize}

\href{https://laravel.com/}{Laravel}\footref{laravelfoot} es un framewok de PHP que se basa en la arquitectura MVC (\ref{MVC}) y tiene como objetivo usar una sintaxis elegante, permitiendo crear código de manera sencilla e implementando muchas funcionalidades\cite{laravel}. Tiene varias utilidades que hemos utilizado dentro del proyecto, por ejemplo las plantillas Blade las cuales permiten usar código PHP e incluir lógica dentro de las vistas de manera muy sencilla y cómoda y también Eloquent ORM el cual nos permite una manera sencilla de manipular los elementos de una base de datos creando un modelo por cada una de las tablas de la base de datos y esto nos permite manipular estos modelos sin tener que ejecutar sentencias SQL. 

\subsection{Gestor de paquetes}

\subsubsection{Composer}

\href{https://getcomposer.org/}{Composer}\footnote{\href{https://getcomposer.org/}{https://getcomposer.org/}} es un gestor de paquetes para PHP, se usa para poder integrar librerías de terceros en tu proyecto de una manera muy sencilla. Para este proyecto se ha utilizado para incorporar las librerías que se han usado de PHP y Laravel (\ref{librerias}) y también el propio Laravel (\ref{Laravel}).

\subsection{Front-end}

\subsubsection{Bootstrap}
\begin{itemize}
    \item Herramientas consideradas: \href{https://getbootstrap.com/}{Bootstrap}\footnote{\label{bootstrapfoot}\href{https://getbootstrap.com/}{https://getbootstrap.com/}}, \href{https://tailwindcss.com/}{Tailwind}\footnote{\href{https://tailwindcss.com/}{https://tailwindcss.com/}}.
	\item Herramienta elegida: \href{https://getbootstrap.com/}{Bootstrap}\footref{bootstrapfoot}.
\end{itemize}

\href{https://getbootstrap.com/}{Bootstrap}\footref{bootstrapfoot} es una librería de HTML, CSS y JavaScript que permite dar estilo a una página web de manera muy sencilla. Esta librería funciona de manera que asignando clases determinadas a los elementos HTML esos elementos va a tener un formato, una posición o un estilo determinado.

\subsubsection{jQuery}

\href{https://jquery.com/}{jQuery}\footnote{\href{https://jquery.com/}{https://jquery.com/}} es una librería de JavaScript la cual nos permite manipular los elementos HTML, añadir animaciones o ejecutar peticiones Ajax de una manera mucho más rápida y sencilla que si lo hiciéramos directamente con Vanilla JavaScript. 

\subsection{Base de datos}

\subsubsection{MySQL}
\begin{itemize}
    \item Herramientas consideradas: \href{https://www.mysql.com/}{MySQL}\footnote{\label{mysqlfoot}\href{https://www.mysql.com/}{https://www.mysql.com/}},
    \href{https://www.postgresql.org/}{PostgreSQL}\footnote{\href{https://www.postgresql.org/}{https://www.postgresql.org/}}.
	\item Herramienta elegida: \href{https://www.mysql.com/}{MySQL}\footref{mysqlfoot}.
\end{itemize}

\href{https://www.mysql.com/}{MySQL}\footref{mysqlfoot} es un sistema gestor de bases de datos relacionales. Se decidió usar MySQL en lugar de PostgreSQL debido al rendimiento de ambos gestores de bases de datos en una aplicación Laravel, como se puede comprobar en la siguiente web (\href{https://medium.com/web-developer/mysql-vs-postgresql-performance-test-with-laravel-api-for-simple-eloquent-queries-on-1-million-6e0e6f1005b8}{Comparación MySQL y PostgreSQL}\footnote{\href{https://medium.com/web-developer/mysql-vs-postgresql-performance-test-with-laravel-api-for-simple-eloquent-queries-on-1-million-6e0e6f1005b8}{https://medium.com/web-developer/mysql-vs-postgresql-performance-test-with-laravel-api-for-simple-eloquent-queries-on-1-million-6e0e6f1005b8}}) MySQL tiene bastante mejor rendimiento.

\subsubsection{PhpMyAdmin}

\href{https://www.phpmyadmin.net/}{PhpMyAdmin}\footnote{\href{https://www.phpmyadmin.net/}{https://www.phpmyadmin.net/}} es una herramienta de gestión de bases de datos MySQL escrita en PHP. Permite ejecutar sentencias SQL mediante una interfaz gráfica muy intuitiva y sencilla de utilizar.


\subsubsection{MySQL Workbench}

\href{https://www.phpmyadmin.net/}{MySQL Workbench}\footnote{\href{https://www.mysql.com/products/workbench/}{https://www.mysql.com/products/workbench/}} es la herramienta oficial de gestión bases MySQL, es similar a PhpMyAdmin a diferencia que esta herramienta tiene una aplicación de escritorio. En principio no se tenía pensado trabajar con esta herramienta, pero a la hora conectar una base de datos a la aplicación web subida en Heroku era necesario usar MySQL Workbench para administrar esta base de datos.

\subsection{Documentación}

\subsubsection{\LaTeX}
\begin{itemize}
    \item Herramientas consideradas: \href{https://www.openoffice.org/es/}{Open Office}\footnote{\href{https://www.openoffice.org/es/}{https://www.openoffice.org/es/}}, \href{https://www.latex-project.org/}{\LaTeX}\footnote{\label{latexfoot}\href{https://www.latex-project.org/}{https://www.latex-project.org/}}.
	\item Herramienta elegida: \href{https://www.latex-project.org/}{\LaTeX}\footref{latexfoot}.
\end{itemize}

\href{https://www.latex-project.org/}{\LaTeX}\footref{latexfoot} es un sistema de composición de textos de alta calidad y su principal uso es para artículos científicos.

Se decidió usar \LaTeX{} por encima de Open Office debido a que se vio como una oportunidad de aprender una herramienta nueva que podía llegar a ser útil para futuros trabajos.

\subsection{Editores de texto}

\subsubsection{Overleaf}

\href{https://es.overleaf.com/}{Overleaf}\footnote{\href{https://es.overleaf.com/}{https://es.overleaf.com/}} es un editor online de \LaTeX{} que permite ir compilando el documento y ver lo cambios al instante, lo cual agiliza mucho el trabajo a la hora de redactar documentos.

\subsubsection{Sublime Text 3}

\begin{itemize}
    \item Herramientas consideradas: \href{https://code.visualstudio.com/}{Visual Studio Code}\footnote{\href{https://code.visualstudio.com/}{https://code.visualstudio.com/}}, \href{https://www.sublimetext.com/3}{Sublime Text 3}\footnote{\label{sublimefoot}\href{https://www.sublimetext.com/3}{https://www.sublimetext.com/3}}.
	\item Herramienta elegida: \href{https://www.sublimetext.com/3}{Sublime Text 3}\footref{sublimefoot}.
\end{itemize}

\href{https://www.sublimetext.com/3}{Sublime Text 3}\footref{sublimefoot} es un editor de texto cómodo, sencillo y fácil de usar, viene incluido con la herramienta Kite que sirve para autocompletar código en distintos lenguajes. Para este proyecto se ha utilizado este editor para redactar todo el código del proyecto, el cual incluye la aplicación web en Laravel, el script Python para hacer la base de datos sintética y los script SQL. Se ha elegido Sublime Text 3 sobré Visual Studio Code debido a que ya había estado usando este editor de texto antes y por lo tanto estaba más familiarizado con el. 

\subsection{Gestión del proyecto}

\subsubsection{GitHub}

\href{https://github.com/}{GitHub}\footnote{\href{https://github.com/}{https://github.com/}} es una herramienta que sirve para subir, almacenar código en la nube y también llevar un control de versiones. GitHub se ha utilizado para llevar a cabo la metodología Scrum gracias a las milestones las cuales he usado para simular los sprints y las issues las cuales son todas las tareas que se han ido realizando a lo largo del proyecto. 

\subsubsection{ZenHub}

\href{https://www.zenhub.com/}{ZenHub}\footnote{\href{https://www.zenhub.com/}{https://www.zenhub.com/}} es una herramienta de gestión de proyecto la cual se integra en GitHub. ZenHub nos permite organizar las issues en pipelines y establecer cuanto tiempo se va a tardar en realizar cada una de las tareas.

\subsection{Análisis de la calidad del código}

\subsubsection{Codacy} \label{Codacy}

\href{https://www.codacy.com}{Codacy}\footnote{\href{https://www.codacy.com}{https://www.codacy.com}} es una herramienta de análisis de código que permite detectar problemas en el código como problemas de seguridad, código repetido o código sin usar.

\subsection{Librerías}\label{librerias}

\subsubsection{Numpy}

Librería de Python que permite crear vectores y matrices y muchas funciones para trabajar con estos. Se ha utilizado en el proyecto para la parte estadística de la creación de la base de datos sintética.  

\subsubsection{mysql.connector}

Librería de Python que permite conectar una base de datos MySQL y realizar operaciones SQL sobre esta Se ha utilizado en el proyecto para insertar los datos en la base de datos sintética.

\subsubsection{phpunit}

Framework de PHP que permite crear pruebas unitarias. Se ha utilizado para realizar las pruebas unitarias de la aplicación web.

\subsubsection{faker}

Librería de Python que permite crear nombre, apellidos y fechas aleatorias. Se ha utilizado para la creación de la base de datos sintética. 

\subsubsection{laravel/ui}

Paquete de autenticación de Laravel. Se ha utilizado para el login de la aplicación web.

\subsubsection{maatwebsite/excel}

Librería de Laravel que permite exportar e importar datos de un fichero Excel. se ha utilizado para poder exportar la base de datos a un fichero en Excel.

\subsubsection{openssl} \label{openssl}

Librería de PHP que permite la encriptación de datos. Se ha utilizado para encriptar los datos sensibles de los pacientes.

\subsubsection{Google Charts}
\begin{itemize}
    \item Herramientas consideradas: \href{https://www.chartjs.org/}{Chart.js}, \href{https://developers.google.com/chart}{Google Charts}.
	\item Herramienta elegida: \href{https://developers.google.com/chart}{Google Charts}.
\end{itemize}

Librería de JavaScript que permite la visualización de datos en gráficas interactivas. Se he utilizado para todo el apartado de visualización de datos.

\subsubsection{DataTables}

Librería de JavaScript que permite la creación de tablas interactivas. Se ha utilizado para todas las tablas que hay implementadas en la web.

\subsubsection{pandas}

Librería de python que permite una sencilla manipulación de conjuntos de datos. Se ha utilizado para el procesamiento de los datos necesarios para dibujar las gráficas kaplan meier.

\subsubsection{lifelines}

Librería de python que permite realizar distintos analisis de supervivencia. Se ha utilizar para dibujar las graficas kaplan meier.




\capitulo{5}{Aspectos relevantes del desarrollo del proyecto}

\section{Reuniones con los oncólogos y preparación del proyecto}

Al principio del proyecto no se conocía muy bien cuál iba a ser el alcance de este ya que solo se sabía el propósito general, una aplicación web que permitiera la gestión de los datos de los pacientes con cáncer de pulmón y que facilitará una mejor visualización de estos. Pero como el desarrollo web era un campo que me llamaba mucho la atención y el hecho que mi proyecto fuera a tener una utilidad real y encima en un campo como la medicina que cualquier mejora o avance puede suponer mejorar la calidad de vida de muchas personas, fue motivo suficiente para elegir este trabajo.

Las dos primeras semanas de proyecto se basaron única y exclusivamente en realizar reuniones con los oncólogos Raquel Gómez Bravo y Enrique Lastra Aras para recoger las historias de usuario del proyecto, entender los datos que se iban a recoger y valorar las distintas posibilidades que había para crear la estructura de la base de datos relacional, adaptar el diseño a lo que ellos buscaban y elegir las herramientas que me fueran a permitir llevar a cabo este proyecto.

Las primeras reuniones fueron complicadas de afrontar debido a que era la primera vez que se trataba con clientes reales y la primera vez que aplicaba los conocimientos adquiridos de ingeniería de software en un escenario real, pero gracias a la colaboración de los oncólogos y a la ayuda de los tutores pudimos acabar entendiéndonos perfectamente y redactando lo que sería el comienzo del proyecto. 

\section{Migraciones, seeders y Scripts SQL}

Al comienzo del desarrollo de la aplicación web se empezó a trabajar en local con el entorno de desarrollo WAMP (\ref{WAMP}), pero con la perspectiva de que en el futuro este proyecto iba a tener que ser integrado tanto en un host online como probablemente en el servidor del HUBU. Por lo tanto no solo iba a tener que realizar la aplicación si no también una manera de poder migrar tanto esta como la base de datos. Esta fue una de las razones por las cuales se selecciono Laravel (\ref{Laravel}) ya que se conocía que tenía dos herramientas para realizar esto de manera sencilla.

La primera herramienta son las migraciones, Laravel nos permite crear clases que hereden de \textit{Migration} con dos funciones, la función \textit{up} la cual sirve para crear la tabla con sus correspondientes campos y referencias (\ref{fig:migracion_up}) y la función \textit{down} que elimina la tabla. Para ejecutar todas las migraciones y crear la estructuras de tablas completa solo hace falta ejecutar el comando \textit{php artisan migrate} en la dirección en la que este instalado nuestro proyecto Laravel \cite{migraciones}.

\imagen{migracion_up}{Método up de la migración de la tabla metástasis}

La segunda herramienta son los \textit{seeders}, Laravel nos permite insertar datos en las tablas ya creadas creando clases que hereden de \textit{Seeder}, estas clases tendrán un solo método llamado \textit{run} el cual insertará los datos dentro de la tabla especificada (\ref{fig:seeders}). Para ejecutar todos los \textit{seeders} y introducir los datos solo hace falta ejecutar el comando \textit{php artisan db:seed}  en la misma dirección que las migraciones \cite{seeders}. 

\imagen{seeders}{Método \textit{run} del \textit{seeder} de los roles}

Por último también se han realizado scripts en SQL para realizar lo mismo que he comentado anteriormente en las migraciones y en los \textit{seeders}, el motivo de esto es para que en el caso que estas herramientas fallen se pueda seguir poniendo a punto la base de datos ejecutando los scripts SQL.

\section{Gestión de errores}

Para este trabajo se han tenido muy en cuenta dos factores, el primero que los médicos que vayan a usar la aplicación es muy probable que no estén familiarizados con la informática y el segundo que a la hora de realizar las representaciones estadísticas tener dos datos iguales pero escritos de manera diferente podía suponer un problema ya que se iban a tomar como datos diferentes.

Para solventar este problema se ha intentado que la mayoría de los formularios de entrada de datos sean desplegables, así eliminando el error humano por completo, y en los campos que esto no era posible, añadir validadores que implementa Laravel para poder comprobar que estos datos están bien introducidos, por ejemplo que la fecha del diagnostico no pueda ser anterior a la fecha de la primera consulta. 

\section{Problemas con la seguridad de los datos sensibles de los pacientes y solución} \label{seguridad}

Una vez completada la primera parte del proyecto, la cual supuso todo el desarrollo de la gestión de usuarios y de pacientes, se pensó que la mejor manera de que los oncólogos fueran usando la aplicación web y buscando tanto errores como posibles mejoras, era subir esta a un host online. Esto en principio no resulto ningún problema, hasta que nos dimos cuenta que podíamos estar vulnerando la ley de privacidad de datos al publicar datos de pacientes sin su consentimiento en servidores externos al hospital. Para solventar este problema se tomaron tres medidas de seguridad.

La primera fue buscar un host online cuyos servidores estuvieran en Europa para que siguiera la misma ley de protección de datos que en España y que este tuviera una política de protección de datos de acuerdo a lo que nosotros necesitábamos. Encontramos Heroku (\ref{Heroku}), el cual es un host de páginas web en la nube con servidores localizados en Irlanda \cite{herokulocalizacion} y la política de privacidad de los datos de esta web dice que su personal no puede acceder ni interactuar con los datos alojados en sus servidores a no ser que se pida permiso previo al usuario o por mandato del gobierno \cite{herokuprivacidad}. Por lo tanto se selecciono este host ya que se vió que cumplía a la perfección con los requisitos de privacidad que teníamos. 

La segunda medida fue cifrar los datos sensibles mediante una encriptación simétrica \ref{simetrica}. Para llevar a cabo esta tarea se usó la librería openssl \ref{openssl} y se creó una clase con 2 métodos, el primero se encarga de encriptar los datos cuando estos son introducidos o modificados en la base de datos y el segundo se encargar de desencriptar los datos cuando estos van a ser visualizados por el usuario en la web. 

La tercera medida fue el protocolo HTTPS ya que la única brecha de seguridad que quedaba en la aplicación era que al usar el protocolo HTTP se podía acceder e interceptar las peticiones que se realizaban desde el cliente al servidor y por lo tanto tener acceso tanto a las contraseñas de los usuarios como a los datos sensibles de los pacientes. Para esto se consiguió un certificado gratuito SSL \cite{ssl}. Pero después de conseguir este certificado apareció otro problema, la versión gratuita de Heroku no permitía integrar un certificado SSL externo y por lo tanto tuve que pasar la cuenta a la versión hobby por un precio de 7\$ para poder integrar el protocolo HTTPS dentro de la web.

\section{Investigación sobre el apartado visual de las gráficas y tablas}

Una de las incertidumbres que surgió a la hora de empezar con el apartado de la visualización de los datos era que no sabía muy bien ni como realizar la interfaz ni como representar los datos obtenidos. Para lo primero, después de darle vueltas se decidió realizar un diseño sencillo el cual permitiera las divisiones según el tipo de dato como se puede ver en la imagen \ref{fig:graficas_1}.  

\imagenTam{graficas_1}{Interfaz para la selección de divisiones para la realización de gráficas}{1}

Para la visualización de las gráficas se investigó sobre diferentes programas estadísticos y al final se decidió por realizar las tablas y gráficas similares al programa \href{https://www.stata.com/}{STATA}\footnote{\href{https://www.stata.com/}{https://www.stata.com/}}.

Como se puede observar en las siguientes imágenes, tenemos por un lado tanto el diagrama de sectores \ref{fig:graficas_2} como la tabla de frecuencias de STATA \ref{fig:graficas_3} y el diagrama de sectores \ref{fig:grafica_normal} y la tabla de frecuencias de nuestra aplicación web \ref{fig:tabla_normal}.

\imagenTam{graficas_2}{Diagrama de sectores STATA}{0.6}
\imagenTam{graficas_3}{Tabla de frecuencias STATA}{0.7}
\imagenTam{grafica_normal}{Diagrama de sectores del proyecto}{0.8}
\imagenTam{tabla_normal}{Tabla de frecuencias del proyecto}{0.8}

Para los atributos numéricos se pidió otro tipo de gráficas y tablas para mejorar el análisis, ambas se basan en los percentiles (\ref{percentil}) y los conceptos teoricos explicado en el apartado \ref{estadistica}. En las siguiente imagenes podemos observar tanto la gráfica de vela \ref{fig:graficas_4} como la tabla de percentiles de STATA \ref{fig:graficas_5} y la gráfica de vela \ref{fig:graficas_percentil} y la tabla de percentiles de nuestra aplicación web \ref{fig:tabla_percentiles}.

\imagenTam{graficas_4}{Gráfica de vela STATA}{0.6}
\imagenTam{graficas_5}{Tabla de percentiles STATA}{0.9}
\imagenTam{graficas_percentil}{Gráfica de vela del proyecto}{0.9}
\imagenTam{tabla_percentiles}{Tabla de percentiles del proyecto}{0.9}

\section{Pruebas}

La realización de pruebas es una parte fundamental en el desarrollo de software para conseguir un producto de calidad y sin errores. Debido a que el objetivo de este proyecto era desarrollar una aplicación web que iba a ser usada por los médicos para la introducción de datos se vió como algo necesario la realización de pruebas que nos aseguren que la web funciona al 100\%, que los datos se almacenan correctamente en la base de datos y que no existe ningún error.

Para esto se llevaron a cabo tres tipos de pruebas: pruebas de integración, pruebas de validación y pruebas de seguridad.

\subsection{Pruebas de integración}

Las pruebas de integración son aquellas que se encargan de probar que los elementos unitarios de código funcionan correctamente en conjunto \cite{pruebasIntegracion}. En nuestro proyecto se han realizado este tipo de pruebas para comprobar la introducción, modificación y eliminación tanto de los datos de los pacientes como de los datos de los usuarios. Para hacer estas comprobaciones se han realizado 138 test los cuales se pueden ejecutar accediendo a la carpeta del proyecto y ejecutando el comando \textit{php artisan test}, en la imagen \ref{fig:test_integracion} podemos ver el resultado de los test. 

\imagen{test_integracion}{Resultado de la ejecución de los test de integración}

\subsection{Pruebas de validación}

Las pruebas de validación son aquellas que se encargan de comprobar que los requisitos y las expectativas del cliente se han cumplido. Este proceso suele realizarse en presencia del cliente o de los usuarios finales \cite{validacion}. Para la validación de este proyecto se han llevado a cabo dos tipos de pruebas: alfa y beta.

\subsubsection{Pruebas alfa}

Las pruebas alfa son aquellas que se realizan en presencia del cliente. En mi caso este tipo de pruebas se realizaron desde mi portátil ejecutando la aplicación en local y en presencia de los oncólogos, a los cuales se les permitió la utilización de la aplicación web con todas sus funcionalidades para comprobar que era exactamente lo que querían. Esta prueba fue mucho mejor de los esperado ya que los únicos errores que se encontraron fueron problemas de nomenclatura en los datos y algún elemento fuera de lugar, es decir problemas que no me iban a llevar mucho tiempo solucionar.

\subsubsection{Pruebas beta}

Las pruebas beta son aquellas realizadas por los usuarios finales y en los equipos de los clientes sin la presencia del desarrollador. En mi caso este tipo de pruebas se realizaron cuando la aplicación web se subió al host online Heroku \ref{Heroku} y los oncólogos, que en el caso de este proyecto eran tanto los clientes como los usuarios finales, pudieron probar la aplicación web al completo en un entorno real.

\subsection{Pruebas de seguridad}

Las pruebas de seguridad son aquellas que tratan de encontrar vulnerabilidades dentro del software. Como ya se comentó en el punto \ref{seguridad}, se había realizado todo lo posible para conseguir una buena seguridad dentro de la aplicación, pero luego me di cuenta que podía existir otro problema y esté era que los middleware no estuvieran funcionando correctamente en alguna de las rutas y cualquier usuario pudiera acceder a ellas sin estar autenticado o sin tener el rol necesario y por lo tanto tener acceso a datos sensibles. Para esto se redactaron pruebas de seguridad en las cuales se realizaron diversas comprobaciones:
\begin{itemize}
    \item Que un usuario sin autenticarse no pueda acceder a ninguna de las rutas de la web.
    \item Que un usuario con rol oncólogo pueda acceder a todas las rutas excepto a las reservadas para administradores.
    \item Que un usuario con el rol administrador tenga acceso a todas las rutas de la web.
\end{itemize}

\textcolor{red}{Estas pruebas se están implementando actualmente.}

\section{Fase de refactorización del código}

\textcolor{red}{Este apartado se completará cuando se acabe la refactorización la cual esta actualmente en proceso.}

\section{Opinión de los oncólogos respecto a la aplicación final y al trabajo realizado}

\textcolor{red}{El miércoles tengo reunión con ellos y en el caso que me deis el visto bueno les pediré que redacten una valoración del trabajo realizado.}


\capitulo{6}{Trabajos relacionados}

Este apartado sería parecido a un estado del arte de una tesis o tesina. En un trabajo final grado no parece obligada su presencia, aunque se puede dejar a juicio del tutor el incluir un pequeño resumen comentado de los trabajos y proyectos ya realizados en el campo del proyecto en curso. 

\capitulo{7}{Conclusiones y Líneas de trabajo futuras}

\subsection{Conclusiones}

Tras haber finalizado el proyecto se puede decir que se han completado todos los objetivos que se habían establecido al principio de este. Como resultado se ha obtenido una aplicación web que permite un eficiente almacenamiento y una sencilla visualización de los datos de pacientes. La idea de este trabajo es ser usado en el HUBU pero con alguna modificación podría implementarse en otros hospitales que tengan el mismo problema.

A nivel personal la realización de este trabajo, sobre todo al principio, ha sido un poco agobiante debido a que se contaba con la presión extra de estar trabajando para unas personas que confiaban en que el proyecto saliera adelante sin ningún problema. Pero a medida que iba avanzando y sobre todo al finalizar el proyecto, el resultado ha sido muy gratificante, ya que no solo he sido capaz de trabajar con esa presión añadida sino que el resultado ha sido bueno y los oncólogos están muy contentos con la labor realizada. También ser consciente que mi aplicación va a ser usada y va ayudar a gente es lo que más positivamente valoro de este proyecto.

A nivel formativo he adquirido conocimientos tanto de diseño como de desarrollo web los cuales son unos campos que desconocía hasta el momento y de los cuales estaba muy interesado en aprender. También al tener pensado seguir llevando el mantenimiento de esta aplicación web se podrá seguir aprendiendo y aumentando los conocimientos sobre estas ramas de la informática.

Como conclusión final puedo decir que estoy bastante contento y orgulloso del trabajo que he llevado a cabo, pero me han quedado algunas cosas pendientes como la inteligencia artificial, que por falta de datos (ya que los oncólogos están empezando ahora a introducirlos en la base de datos) no se ha podido llevar a cabo y también me hubiera gustado disponer de más tiempo para pulir más el diseño y añadir algún tipo de gráfica más. 

\subsection{Lineas de trabajo futuras}

Con los oncólogos del HUBU se acordaron unas lineas de trabajo futuras ya que por contratiempos y por temas burocráticos ha habido tareas que han quedado aún pendientes. A continuación voy a listar todas estas tareas:

\begin{itemize}
    \item Migrar tanto los datos ya introducidos por los oncólogos en la base de datos como la aplicación web al servidor del HUBU. 
    \item Añadir nuevas gráficas que los oncólogos vean necesarias en el futuro.
    \item Continuar con el mantenimiento de la web e implementando todas aquellas mejoras o funcionalidades que los oncólogos necesiten.
    \item Mejorar el diseño de la web haciéndola más intuitiva y adaptándola a las necesidades de los oncólogos.
    \item Realizar algún modelo de predicción el cual se entrene con los datos introducidos en la aplicación por lo oncólogos y permita recomendar los tratamientos o medicamentos a cada paciente según sus características.
\end{itemize}



\bibliographystyle{plain}
\bibliography{bibliografia}

\end{document}
